При \(w=0{,}40\) на 1242~К оба корня внутри интервала приложенных полей найдены в 61 из 64 блочных повторов. Оценки ширины равны \(+0{,}446\)~п.п. [0,301; 0,541] на Ti-богатой ветви и \(-0{,}380\)~п.п. [\(-0{,}439\); \(-0{,}276\)] на Al-богатой. На 1244~К максимум \(S=-0{,}605\) [\(-1{,}138\); \(-0{,}071\)], и в центральной оценке ни одного корня нет. На этой изотерме оба корня появились лишь в одном блочном повторе.

На 1243~К обе стенки дали оба корня в 64 из 64 повторов объединённой модели. При \(w=0{,}35\) поля контактов равны 3,47251 и 3,61749~эВ. На Ti-богатой стороне два соседних независимых расчёта имеют противоположные знаки даже с учётом 95\%-х интервалов. На Al-богатой стороне знаки центральных оценок также противоположны, но \(D(3{,}624)=-0{,}45\) [\(-1{,}04\); 0,34]. При \(w=0{,}40\) поля 3,52767 и 3,56297~эВ разрешает объединённая девятиполевая сетка; отдельные соседние пары сами по себе слабее подтверждают смену знака. На 1262~К при \(w=0{,}35\) оба центральных корня присутствуют, однако внутри сетки они сохранились только в 45 из 64 блочных повторов.

При поочерёдном исключении любого внутреннего поля обе границы на 1243~К сохраняются. Для узкой стенки ширины Ti-/Al-богатой ветви меняются в пределах \([+1{,}30;\,+1{,}35]\) и \([-0{,}76;\,-0{,}74]\)~п.п., для широкой --- \([+0{,}24;\,+0{,}30]\) и \([-0{,}25;\,-0{,}21]\)~п.п.